\begin{phasebox}
\textbf{هدف:} اجرای تکرارپذیر پروژه در محیط محلی و ثبت قدم بعدی.\\
\textbf{پیش‌نیاز:} تست‌های نسخه اولیه موفق باشند.\\
\textbf{معیار پایان:} پروژه از یک نسخه تازه مخزن راه‌اندازی و ماندگاری داده بررسی شود.
\end{phasebox}

\section{فایل‌های زمان اجرا}
پایگاه داده، رخدادها و اطلاعات موقت در پوشه
\lr{\texttt{runtime}}
قرار می‌گیرند. این فایل‌ها وارد
\lr{Git}
نمی‌شوند. دستورهای راه‌اندازی در فایل راهنمای هر پوشه نوشته خواهند شد.

\section{آزمایش راه‌اندازی}
یک نسخه تازه از مخزن دریافت می‌شود. سپس تنظیمات محلی ساخته می‌شوند، اولین پیام دریافت می‌شود، داده در وب نمایش داده می‌شود و یک فرمان آزمایشی ارسال می‌گردد. پس از اجرای دوباره برنامه، ماندگاری داده بررسی می‌شود.

\section{پشتیبان‌گیری}
روش پشتیبان‌گیری و بازیابی فایل
\lr{SQLite}
قبل از استفاده طولانی‌مدت ثبت و یک بار آزمایش می‌شود.

\section{توسعه آینده}
\begin{itemize}
    \item برای افزودن برد جدید، یک پوشه هم‌سطح با پوشه آردوینو ساخته می‌شود.
    \item برای تغییر پایگاه داده، کد پوشه پایگاه داده جایگزین می‌شود.
    \item برای تغییر
    \lr{Flask}
    به
    \lr{Django}
    فقط پوشه برنامه وب بازنویسی می‌شود.
    \item تحلیل داده، به‌روزرسانی بی‌سیم و استقرار ابری پس از ایجاد نیاز واقعی اضافه می‌شوند.
\end{itemize}

\section{ثبت نتیجه نهایی}
نسخه برنامه، نسخه کد برد، محدودیت‌های شناخته‌شده، نتیجه تست‌ها و قدم بعدی دقیق در این بخش ثبت می‌شوند.
